\scanpage{115}
\deleted{En fait, on demande qu'} On applique $\tau_\infty=\mathrm R^*\tau_\infty$, \unsure{on tire}
\[
(2)\qquad\diagram{\begin{tikzcd}[column sep=large,row sep=large]
\deleted{\mathrm R^*}\tau_\infty(\mathrm R^*\Gamma(X,S_k))\arrow[r]&\unsure{\tau_\infty}\mathrm R^*\Gamma(X,C^{**})\\
&\tau_\infty C^{**}(X)\arrow[u]
\end{tikzcd}}
\]
On \deleted{demande} \inserted{voudrait} que ces deux flèches soient des isomorphismes \inserted{dans la cat. dérivée}, ce qui signifie aussi que les hom. déduits de (1) en passant à la \unsure{coh.}
\[
(1')\qquad\diagram{\begin{tikzcd}[column sep=large,row sep=large]
H^{**}(X,S_k)\arrow[r]&H^{**}(X,C^{**})\\
&\mathcal H_C^{**}(X)\arrow[u]
\end{tikzcd}}
\]
soient à noyau et conoyau négligeables, ce qui \deleted{signifie} \inserted{implique} précisément, pour la flèche horizontale, l'iso.
\[
\unsure{(2')}\quad\left\{
\diagram{\begin{tikzcd}[column sep=large,row sep=large]
\tau_\infty(H^*(X,k)\otimes_k S_k)\arrow[r,"\sim"]\arrow[dr,dashed,no head]&\tau_\infty H^{**}(X,C^{**})\\
&\tau_\infty(\mathcal H_C^{**}(X))\arrow[u,"\sim"']
\end{tikzcd}}
\right.
\]
\margin{114}
